\section{はじめに}
\label{sec_intro}
大喜利は，出題された題目に対して機知に富んだ回答を返す言葉遊びである．
大喜利投稿サイト「ボケて」\cite{ref_bokete}では，投稿者がお題となる画像を掲示し，他の利用者がその画像に対する面白い一言を投稿する．
以下では，お題となる画像をお題画像，利用者が投稿する一言をボケと呼ぶ．
ボケてに蓄積されたお題画像とボケの組は，画像を入力として面白い文を生成するモデルの学習や，ユーモアの分析に利用できる大規模な資源である．
これらの用途では，各お題画像がどの種類に属するかを示すラベルが必要になる．
本研究では，約682万件のお題画像にラベルを付与する方法を構築してきた．

お題画像の扱いを最も難しくしている要素は，画像中の文字である．
お題画像の約6割には読める文字がある．
その文字には，投稿者が回答者に向けて書いた問いや指示もあれば，写真に写った看板，テレビの字幕，漫画の台詞のように，素材に元からある文字もある．
回答者にとって前者は答えるべき題であり，後者は一言を付ける対象の一部である．
したがって，ボケの生成では，文字を読み取るかだけでなく，読み取った文字を題として扱うかを決める必要がある．

従来のラベルでは，題として扱う文字を「お題文」と呼び，投稿者が回答者に向けて書き加えた文と定義していた．
しかし，この定義は投稿者の意図に依存する．
例えば，アニメの画像に重ねた問いの字幕は，元の放送の字幕である場合も，投稿者がテロップ風に書き加えた場合もあり，両者は見た目では区別できない．
このような判断の揺れは，正解のラベルの再現性を下げるだけでなく，教師となる視覚言語モデル (Vision-Language Model，VLM) の判定や，その判定を学習する生徒の分類器の精度も下げる．

本稿では，お題画像の文字の扱いについて，定義の変遷と，各段階の定義がラベルの付与に与えた影響をまとめる．
本稿の貢献は次のとおりである．
\begin{itemize}
    \item ラベルの定義の改訂の過程で，文字の扱いをどのように変えたかを整理し，投稿者の意図に基づく定義の下で生じた誤りの傾向を示す．
    \item 文字を，写っている場所による層と，文の形による形の2軸の属性に分け，お題文の有無を属性から規則で導く定義を提案する．
    \item 提案の属性の再現性，従来の定義との差，規則の選択による全データのラベルの分布の変化を定量的に示す．
    \item 提案の属性を VLM の教師と画像の埋め込みや書き起こしに基づく生徒で学べるかを調べ，移行の障害を明らかにする．
\end{itemize}

本稿の構成は次のとおりである．
\ref{sec_data}章で対象のデータと用語を定義し，\ref{sec_related}章で関連研究を整理する．
\ref{sec_history}章で文字の扱いの変遷を述べ，\ref{sec_intent}章で投稿者の意図に基づく定義の影響を分析する．
\ref{sec_attr}章で文の属性による定義を提案し，\ref{sec_exp}章と\ref{sec_results}章でその評価の方法と結果を示す．
\ref{sec_discussion}章で考察し，\ref{sec_conclusion}章でまとめる．
